\scanpage{124}
\mspage{8}
% archive sheet number: 123
en un sens, le cas général doit se ramener à ce cas particulier banal.

Le seul type \unsure{d'autres|d'autre} cas que je connaisse, est celui où $\unsure{\mathscr H|\mathcal H|H}=MW^{-1}$, où $W$ est un localiseur admettant une théorie de la fibre $W$-homotopique (cf $\underline{\overline{\mathrm{XII}}}$). Dans ce cas, on l'a dit, on prend comme carrés $h$-cartésiens dans $\mathscr H$ ceux qui sont isomorphes, dans $\underline{\operatorname{Hom}}(\unsure{\mathbb{\Delta}|\Delta}^{1}\times\mathbb{\Delta}^{1},\mathscr H)$, à un carré $W$-cartésien dans $M$. Encore faut-il vérifier les axiomes $\operatorname{Cart}(1)$ à $\operatorname{Cart}(5)$. Je pense que je n'ai guère de chance d'y arriver que moyennant une

\uline{Proposition 1}. Moyennant des hypothèses sur $(M,W)$ \unsure{pouvant|paraissant|précisant} être dégagées ci-dessous, vérifiées en tous cas où $(M,W)$ est \unsure{quarrinable|quarrable|quarrifiable} \deleted{\unsure{dég. de|degré|degré de} \unsure{uniformément|uniformité}} sur la catégorie des ens. ordonnés finis \inserted{(cf $\underline{\overline{\mathrm{XIII}}}$, p. 51)}, on a ceci :
% vertical emphasis rule alongside the proposition
\margin{$\vert$}
soit
\diagram{\begin{tikzcd}[ampersand replacement=\&,execute at end picture={\node[anchor=east] at ([xshift=-5pt]current bounding box.west) {$X=\Biggl($};\node[anchor=west] at ([xshift=5pt]current bounding box.east) {$\Biggr)$};}]
X_1 \arrow[d] \& X_3 \arrow[l] \arrow[d] \\
X_0 \& X_2 \arrow[l]
\end{tikzcd}}
un carré dans $M$, et soit $\mathrm{ho}_W(X)$ son image dans $\underline{\operatorname{Hom}}(\mathbb{\Delta}^{1}\times\mathbb{\Delta}^{1},\mathrm{Ho}_W)$.
% second vertical emphasis rule alongside the concluding lines
\margin{$\vert$}
Pour que $\mathrm{ho}_W(X)$ soit $h$-cartésien, il \uline{faut} (et il suffit, bien sûr) que $X$ soit $W$-cartésien.
